\section{Conclusion}
\label{sec:conclusion}

We have introduced \sa{} bisection, a Structure-layer algorithm that replaces
the single \metis{} call at each bisection tree node with a simulated
annealing refinement loop.
The key design decisions---\metis{} initialisation, geometric cooling scaled
to the subgraph size, and best-ever plan tracking---combine to give a
practical optimiser that achieves 10--18\% Polsby-Popper improvement over
baseline \metis{} at approximately $12\times$ the runtime cost.

\paragraph{Position in the B-series algorithm landscape.}
\sa{} is a Structure-layer optimiser, analogous to \geosec{} and \areasec{}
but targeting the solver rather than the objective or tree structure.
It is orthogonal to SeedCompositor strategies (\cs{}, \be{}) and can be
composed with them.
The empirical quality ordering is:
\[
  \text{METIS} \;<\; \text{SA(steps=10)} \;<\; \text{SA(steps=50)}
  \;<\; \text{BE}(T=100),
\]
with runtimes following the same order.
For practitioners, this ordering suggests a practical decision rule:
\begin{itemize}
\item \textbf{Fastest}: baseline \metis{} (4 s for NC).
\item \textbf{Single-pass best}: SA(steps=10) (52 s for NC), recommended for
  publication-quality single-plan generation.
\item \textbf{Maximum quality, moderate cost}: \be{}($T=100$) (420 s for NC).
\item \textbf{Maximum quality, high cost}: \cs{} or \be{} + SA (hours).
\end{itemize}

\paragraph{Optimiser vs.\ sampler.}
SA is a deterministic optimiser: it returns the best single bisection found
in $n_{\mathrm{steps}}$ steps from a fixed seed.
\recom{} is a Markov chain sampler: it generates a distribution over valid
plans.
These tools answer different questions.
Use \sa{} (or \be{}) when the question is ``what is the best achievable
redistricting plan under this objective?''
Use \recom{} (or ensemble methods) when the question is ``what does the
space of valid redistricting plans look like, and where does the enacted
map fall?''

\paragraph{Implementation.}
\sa{} bisection is implemented in the \texttt{redist} binary under
\texttt{--structure simulated-annealing}.
The SA parameters are exposed as \texttt{--sa-steps-per-tract} (default 10)
and \texttt{--sa-t0-factor} (default 0.01).
The YAML configuration format supports:
\begin{verbatim}
algorithm:
  structure: simulated-annealing
  sa_steps_per_tract: 10
  sa_t0_factor: 0.01
\end{verbatim}

\paragraph{Open questions.}
Several extensions are straightforward to implement and are left for future work:
(1) combining SA bisection with the \ar{} prime-factor tree structure (B.11),
    which links the bisection tree to the Huntington-Hill apportionment and
    may create larger, harder subgraphs at prime-split nodes where SA has
    more room to improve;
(2) a VRA-aware SA objective that incorporates minority population constraints
    directly into the acceptance criterion, rather than treating VRA as a
    post-hoc filter;
(3) parallel SA: running SA independently on multiple cores and selecting the
    best result is equivalent to \be{} with SA bisectors, and could be
    explored as a first-class configuration option.

These extensions inherit the clean three-layer compositor design of the
existing platform and require no structural changes to the codebase.
